%%%%%%%%%%%%%%%%%%%%%%%%%%%%%%%%%%%%%%%%%%%%%%%%%%%%%%%%%%%%%%%%%%%%%%%%%%%%%%%%
% JOURNAL — ENTRY 001
% GER — Geometria Espectral Relacional
%%%%%%%%%%%%%%%%%%%%%%%%%%%%%%%%%%%%%%%%%%%%%%%%%%%%%%%%%%%%%%%%%%%%%%%%%%%%%%%%

\section{Journal Entry 001}
\label{sec:journal001}

\subsection*{Date}
14 July 2026

\subsection*{Series}
S27 --- Mathematical Foundations

\subsection*{Status}
In Progress

\subsection*{Objective}

This journal records the mathematical development of the S27 series.
Unlike the main manuscript, which contains only consolidated results, this document preserves the chronological evolution of definitions, propositions, audits, corrections, and structural decisions.

The purpose of this journal is to provide complete scientific traceability of the development process.

\subsection*{Current Milestone}

Beginning the formal mathematical foundation of the observational framework of GER.

Primary objectives:

\begin{itemize}
    \item Formalize observational mappings as mathematical applications.
    \item Build the induced equivalence relations.
    \item Construct quotient spaces.
    \item Develop the algebra of observational partitions.
    \item Separate universal mathematics from GER-specific modeling choices.
    \item Audit logical consistency and minimality of the framework.
\end{itemize}

\subsection*{Development Log}

%============================================================
\section{S27-E1 --- Experimental Validation of Operator-Induced Partitions}

With the mathematical foundation of the operator-induced partition formalism established in Series S27-A, the first experimental phase was initiated.

The objective of Experiment S27-E1 was not to investigate new physics, but rather to validate the computational implementation of the partition formalism using signatures generated by the validated S26 framework.

The experiment consisted of constructing the equivalence partitions induced independently by each Geometric Fundamental Operator (GFO):

\[
D,\qquad
C,\qquad
R,\qquad
Dr,
\]

where

\[
D=\text{Diameter},
\]

\[
C=\text{Convergence},
\]

\[
R=\text{Recurrence},
\]

and

\[
Dr=\text{Drift}.
\]

For each operator, the partition

\[
\mathcal P_O
\]

was built directly from the set of geometric signatures generated by the S26 Geometry Scan.

A total of

\[
56
\]

independent signatures were analyzed.

The resulting partitions are summarized below.

\begin{center}
\begin{tabular}{lccc}
\hline
Operator & Blocks & Largest Block & Smallest Block\\
\hline
Diameter      & 56 & 1 & 1\\
Convergence   & 56 & 1 & 1\\
Recurrence    & 11 & 35 & 1\\
Drift         & 56 & 1 & 1\\
\hline
\end{tabular}
\end{center}

Three operators (Diameter, Convergence and Drift) produced completely discrete partitions over the analyzed dataset, indicating that every generated signature occupied a distinct equivalence class.

In contrast, the Recurrence operator generated only eleven equivalence classes, one of which contained thirty-five distinct signatures.

This result immediately suggested the existence of strong observational degeneracy associated with Recurrence.

%============================================================
\subsection{S27-E1.2 --- Recurrence Degeneracy Analysis}

To investigate the origin of the observed degeneracy, the internal variability of each recurrence class was measured.

The largest equivalence class corresponded to

\[
R=0.001282051282051282
\]

and contained

\[
35
\]

different signatures.

Inside this single recurrence class, the remaining operators exhibited large variability.

\begin{center}
\begin{tabular}{lc}
\hline
Observable & Range\\
\hline
Diameter &
$1.809466 \rightarrow 131.629819$\\

Convergence &
$185.668083 \rightarrow 13501.195820$\\

Drift &
$0.993518 \rightarrow 1.000000$\\
\hline
\end{tabular}
\end{center}

Therefore, signatures that are observationally identical with respect to Recurrence remain highly distinguishable under Diameter, Convergence and Drift.

This demonstrates that the partition

\[
\mathcal P_R
\]

is substantially coarser than the partitions induced by the remaining operators.

Consequently, Recurrence alone is not sufficient to distinguish the geometric macrostates generated by the GER engine.

This constitutes the first experimental evidence obtained in Series S27 that different Geometric Fundamental Operators possess substantially different observational resolution.

The natural continuation of this investigation is the construction of joint partitions through the meet operation,

\[
\mathcal P_D
\wedge
\mathcal P_C
\wedge
\mathcal P_R
\]

in order to determine whether the observed degeneracy disappears when multiple operators are considered simultaneously.
%============================================================
%==========================================================

\subsection{S27-E1.3 --- Partition Refinement}

The joint partitions

\[
\mathcal P_D
\wedge
\mathcal P_C,
\qquad
\mathcal P_D
\wedge
\mathcal P_C
\wedge
\mathcal P_R,
\qquad
\mathcal P_D
\wedge
\mathcal P_C
\wedge
\mathcal P_R
\wedge
\mathcal P_{Drift}
\]

were constructed in order to verify whether combining multiple Geometric Fundamental Operators increases the observational resolution.

For the complete Geometry Scan dataset, all refined partitions contained

\[
56
\]

equivalence classes, identical to the partitions induced individually by Diameter, Convergence and Drift.

Therefore, no additional observational resolution is obtained by combining these operators.

This result indicates that these three observables induce the same maximal partition over the investigated dataset.

%==========================================================

\subsection{S27-E2 --- Partition Lattice}

The refinement relations between the observational partitions were investigated through the lattice meet operation.

The experiments confirmed that

\[
\mathcal P_D
=
\mathcal P_C
=
\mathcal P_{Drift},
\]

while

\[
\mathcal P_D
\preceq
\mathcal P_R.
\]

Subsequent lattice consistency audits verified reflexivity, antisymmetry and meet consistency for all operators.

The resulting partition lattice satisfied all tested algebraic properties, providing the first internally consistent lattice structure generated by the GER observational operators.

%==========================================================

\subsection{S27-E3 --- Robustness of the Partition Structure}

A series of robustness experiments investigated whether the partition lattice depends on numerical implementation choices.

The analysis included variations of integration time, timestep, parameter expansion and internal pipeline audits.

In every experiment the refinement relations remained unchanged.

Consequently, the partition lattice appears to represent a structural property of the generated signatures rather than a numerical artifact of the simulation procedure.

%==========================================================

%==========================================================

\subsection{S27-E4 --- Emergence of a Resolution-Dependent Geometric Observable}

After establishing the robustness of the partition lattice under changes of numerical resolution, integration timestep, parameter expansion and internal implementation audits, the investigation focused on the Recurrence operator itself.

The recurrence observable is computed through a resolution parameter

\[
\varepsilon
=
0.05\,\sigma(\Gamma),
\]

where $\sigma(\Gamma)$ denotes the standard deviation of the trajectory.

A systematic resolution scan demonstrated that the previously identified zero-recurrence states do not remain invariant under changes of $\varepsilon$.

For all nine signatures initially satisfying

\[
R=0,
\]

the recurrence remained exactly zero for

\[
\varepsilon
=
0.01\sigma,
\;
0.02\sigma,
\;
0.05\sigma,
\]

while positive recurrence values emerged for larger resolutions.

Therefore, the zero-recurrence regime is not an absolute dynamical property but rather a resolution-dependent observational phenomenon.

This observation motivated the introduction of the \emph{critical recurrence resolution},

\[
\varepsilon^\ast,
\]

defined as the smallest observational resolution for which the recurrence becomes strictly positive.

The critical resolution exists for all signatures generated by the Geometry Scan and exhibits substantial variability,

\[
0.010
\le
\varepsilon^\ast
\le
0.097,
\]

with

\[
\langle
\varepsilon^\ast
\rangle
=
0.0315.
\]

Correlation analysis showed that $\varepsilon^\ast$ is only weakly correlated with the previously established observables (Diameter, Convergence, Drift) as well as with the control parameters $\beta$ and $\sigma$, suggesting that it captures structural information not represented by the existing signature components.

Finally, the partition induced by the critical resolution,

\[
\mathcal P_{\varepsilon^\ast},
\]

contains

\[
30
\]

equivalence classes.

This partition occupies an intermediate position between the fine partitions generated by Diameter, Convergence and Drift,

\[
56
\]

classes,

and the coarse partition generated by Recurrence,

\[
11
\]

classes.

Moreover,

\[
\mathcal P_{\varepsilon^\ast}
\]

is not comparable to

\[
\mathcal P_R
\]

under the refinement order, demonstrating that both observables encode distinct structural aspects of the geometric trajectories.

These results provide the first evidence that observational resolution itself may induce an independent geometric descriptor within the GER framework.

%==========================================================
%==========================================================

\subsection{S27-E5 --- Minimal Observational Structure}

The final stage of Series S27 investigated whether all observational operators provide genuinely independent geometric information.

The individual partitions generated by Diameter, Convergence, Drift, Recurrence and the Critical Recurrence Resolution were compared under the partition equivalence relation.

The experiments demonstrated that

\[
\mathcal P_D
=
\mathcal P_C
=
\mathcal P_{Drift},
\]

indicating that these three observables induce exactly the same observational partition over the complete Geometry Scan dataset.

Consequently, they belong to a single observational family despite representing distinct numerical quantities.

The remaining operators,

\[
\mathcal P_R
\]

and

\[
\mathcal P_{\varepsilon^\ast},
\]

generate distinct observational structures.

Therefore, the five investigated observables reduce to three independent observational families,

\[
\{
D,
C,
Drift
\},
\qquad
\{
R
\},
\qquad
\{
\varepsilon^\ast
\}.
\]

This result provides the first experimental decomposition of the GER observational structure into equivalence classes of geometric operators.

Rather than introducing additional independent observables, the analysis reveals that the observational geometry generated by the current GER framework possesses three irreducible observational families.

This concludes the structural investigation initiated in Series S27.

%==========================================================

\subsection{S27-RA --- Architectural Preparation for External Validation}

Following the completion of Series S27, the objective of the project shifted from the internal characterization of the GER framework to its external validation.

The original intention was to directly apply the GER observational pipeline to classical dynamical systems. However, the first implementation attempt immediately revealed an architectural limitation.

The Persistence Observatory does not operate directly on trajectories.

Instead, it receives a spectral snapshot containing the quantities

\[
\{
\gamma,\;
P,\;
PR,\;
C_m
\},
\]

where \(P\) denotes the modal probability distribution, \(PR\) the participation ratio and \(C_m\) the modal center.

This observation motivated a complete architectural audit of the interface between the simulation engine and the observational framework.

The audit demonstrated that the GER observatory is not directly coupled to the relational engine itself, but rather to a modal representation of the system state.

The snapshot consumed by the observatory is generated through a modal projection,

\[
\hat{\gamma}
=
V^{T}\gamma,
\]

where \(V\) denotes the modal basis.

This observation naturally led to the formulation of the Modal Embedding Hypothesis.

Rather than requiring the specific relational Laplacian of the GER engine, the observational framework appears to require only a modal representation satisfying a small set of mathematical properties.

A systematic architectural analysis identified the following requirements:

\begin{itemize}
\item complete state representation;
\item intrinsic ordering of modes;
\item coefficient normalization;
\item construction of a probability distribution;
\item well-defined participation ratio;
\item well-defined modal center;
\item well-defined spectral entropy.
\end{itemize}

Candidate modal representations were then analyzed according to these criteria.

Among the investigated candidates, the Fourier basis satisfied every currently identified mathematical requirement.

Finally, the projection convention employed by the GER framework was experimentally audited.

Using the Fourier basis, the modal projection

\[
V^{T}\gamma
\]

reproduced the discrete Fourier transform with numerical errors at machine precision,

\[
\|\Delta\|
\approx
4.7\times10^{-16},
\]

demonstrating that the projection architecture implemented by the GER framework is fully compatible with a Fourier modal basis.

Consequently, the architectural dependence of the GER observatory was shown to be associated with modal representations rather than with the specific relational dynamics that generated them.

This result does not yet constitute an external validation of the GER framework.

Instead, it establishes the mathematical and architectural foundations required to perform external validation without modifying any component of the observational pipeline.

The first external observation experiments are therefore enabled under a rigorously defined modal embedding protocol.

%==========================================================
\subsection{S27-R1.5 — First External Geometric Signature}

The first complete external execution of the GER observational pipeline was successfully performed.

Unlike every previous experiment in the GER framework, the analyzed system was not generated by the relational engine itself. Instead, an independent Harmonic Oscillator was represented through a compatible modal embedding and processed without any modification of the observational framework.

The complete observational pipeline executed was:

\[
\text{External System}
\rightarrow
\text{State Representation}
\rightarrow
\text{Observational Snapshot}
\rightarrow
\text{Persistence Observatory}
\rightarrow
\text{Observational Trajectory}
\rightarrow
\text{Fundamental Geometric Operators}
\rightarrow
\text{Geometric Signature}.
\]

No system-specific observables, operators or calibration parameters were introduced at any stage of the pipeline.

During Reality Validation, an architectural limitation of the observatory was identified. The modal energy computation implicitly assumed real-valued modal coefficients,

\[
E_i = c_i^2,
\]

which is valid only for real spectral bases. The implementation was generalized to

\[
E_i = |c_i|^2,
\]

allowing the observational framework to operate correctly on complex modal bases such as the Fourier representation while preserving compatibility with all previous analyses based on real eigenvectors.

The Harmonic Oscillator produced the following Geometric Signature:

\[
\begin{aligned}
\text{Diameter} &= 1.9921147013144778,\\
\text{Convergence} &= 0.08300477922143658,\\
\text{Recurrence} &= 0.022108843537414966,\\
\text{Drift} &= 4.062050898015295\times10^{-17}.
\end{aligned}
\]

The nearly vanishing Drift is consistent with the bounded periodic dynamics expected for the harmonic oscillator. Likewise, the low Convergence and finite Recurrence indicate that the trajectory remains confined without exhibiting asymptotic collapse or irreversible evolution.

Most importantly, this experiment demonstrates that the complete observational structure of GER is capable of generating a valid Geometric Signature for a dynamical system that did not participate in the development of the framework.

This result does not validate the GER theory itself. Rather, it constitutes the first external validation of the observational pipeline and establishes that the same Fundamental Geometric Operators can be applied consistently to an independent dynamical system without introducing system-specific modifications.

This experiment therefore marks the completion of the first full external observational cycle of the GER framework, concluding the Geometric Signature stage of the Reality Validation Series (S27-R).

\subsection{S27-R1.6 — First External Structural Certificate}

Following the successful generation of the first external Geometric Signature, the complete Structural Certification pipeline of GER was evaluated using an independent dynamical system.

The Harmonic Oscillator signature generated in Experiment S27-R1.5 was submitted directly to the official Structural Certification module without introducing any system-specific modifications.

The executed pipeline was therefore

\[
\text{External System}
\rightarrow
\text{Modal Representation}
\rightarrow
\text{Observational Snapshot}
\rightarrow
\text{Persistence Observatory}
\rightarrow
\text{Observational Trajectory}
\rightarrow
\text{Fundamental Geometric Operators}
\rightarrow
\text{Geometric Signature}
\rightarrow
\text{Structural Certificate}.
\]

During this experiment, Reality Validation identified an architectural incompatibility between the public \texttt{Signature} API and the Structural Certification module.

The official \texttt{Signature} object had previously been migrated to an immutable dataclass representation, whereas the certification pipeline still assumed an iterable dictionary interface through expressions of the form

\[
\texttt{dict(signature)}.
\]

Rather than modifying the certification algorithms, backward compatibility was restored directly at the API level by extending the \texttt{Signature} interface while preserving its immutable representation.

This correction maintained full compatibility with existing scientific modules without altering the mathematical structure of the certification process.

The external Harmonic Oscillator was then successfully certified.

The resulting Structural Certificate reported

\[
\begin{aligned}
\text{Signature Integrity} &:\ \mathrm{PASS},\\
\text{Consistency Check} &:\ \mathrm{PASS}.
\end{aligned}
\]

The Signature Integrity operator confirmed that the externally generated Geometric Signature satisfies all structural requirements expected by the GER certification framework.

The Consistency Check further verified that the generated certificate is fully compatible with the current formal definition of the Fundamental Geometric Operators (OGFs).

An important consequence follows from this experiment.

The Structural Certification module does not depend on the origin of the analyzed system. Once a compatible Geometric Signature is produced, the certification process operates exclusively on its geometric structure.

Therefore, at least for the Harmonic Oscillator, the GER Structural Certification framework behaves as a system-independent geometric validator rather than a validator restricted to internally generated relational simulations.

This experiment completes the first fully external execution of the complete GER observational and certification pipeline and constitutes the first successful external Structural Certificate produced by the framework.
This result establishes the first externally validated reference point within the observational geometry of GER.

Rather than validating an isolated computational experiment, S27-R1 demonstrates that an independent physical system can successfully traverse the complete GER observational pipeline—from modal representation and spectral observation to geometric characterization and structural certification—without requiring any system-specific modification of the framework.

The successful generation of an external Geometric Signature, its subsequent Structural Certification, and its full compatibility with the official observational infrastructure indicate that the Fundamental Geometric Operators operate exclusively on the intrinsic geometric structure of the observed dynamics rather than on the internal origin of the analyzed system.

Consequently, the completion of S27-R1 marks the first experimentally validated external element of the GER observational universe.

Having established the complete external validation of the observational pipeline, the subsequent stages of the S27-R protocol focus on investigating how progressively more complex dynamical systems populate the same observational geometry and whether distinct physical regimes naturally organize into identifiable geometric structures under the GER framework.

\subsection{S27-R2 — External Validation Using a Damped Oscillator}

Following the successful external validation of the Harmonic Oscillator, the complete GER observational pipeline was evaluated using a damped harmonic oscillator, introducing dissipative dynamics while preserving the same experimental protocol.

No modifications were introduced into the GER framework. The external trajectory was processed through the identical observational chain,

\[
\text{External System}
\rightarrow
\text{Modal Representation}
\rightarrow
\text{Observational Snapshot}
\rightarrow
\text{Persistence Observatory}
\rightarrow
\text{Observational Trajectory}
\rightarrow
\text{Fundamental Geometric Operators}
\rightarrow
\text{Geometric Signature}
\rightarrow
\text{Structural Certificate}
\rightarrow
\text{Observational Partition}.
\]

The Persistence Observatory detected measurable differences with respect to the conservative system, particularly through the local residual observable, indicating that the observational metrics respond to dissipative evolution while remaining numerically stable.

The resulting Geometric Signature was

\[
\begin{aligned}
D &= 1.9998041350423283,\\
C &= 0.6184311298641284,\\
R &= 0.027210884353741496,\\
\Delta &= 0.00807519727975368.
\end{aligned}
\]

Comparison with the Harmonic Oscillator revealed that the Geometric Signature changed consistently under dissipation. In particular, both the Convergence and Drift operators exhibited significant variations, while the Diameter remained nearly unchanged, indicating that the Fundamental Geometric Operators are sensitive to dynamical regime rather than merely to the spatial extent of the observed trajectory.

The generated signature successfully passed the official Structural Certification process, producing

\[
\begin{aligned}
\text{Signature Integrity} &:\ \mathrm{PASS},\\
\text{Consistency Check} &:\ \mathrm{PASS}.
\end{aligned}
\]

As in the previous experiment, certification depended exclusively on the geometric structure of the produced signature and remained completely independent of the physical origin of the analyzed system.

The certified signature was subsequently incorporated into the official observational universe generated by GER and partitioned using the standard observational operators. Under the recurrence-induced partition, the external Damped Oscillator formed a singleton observational block, demonstrating that the observational partition framework consistently accommodates external dissipative dynamics without requiring any modification of the underlying theoretical infrastructure.

The completion of S27-R2 extends the external validation of GER beyond conservative dynamics and provides the first successful observational and structural validation of a dissipative physical system. Together with S27-R1, these experiments demonstrate that the GER observational framework operates consistently across distinct classes of external dynamical systems while preserving a unified geometric and structural interpretation.

\subsection{S27-R3 --- External Validation Using the Van der Pol Oscillator}

The third stage of the external validation protocol evaluated GER using the classical Van der Pol oscillator, introducing the first nonlinear dynamical system into the observational framework.

Unlike the previous experiments, the Van der Pol oscillator exhibits self-sustained oscillations generated by a nonlinear restoring mechanism, converging asymptotically to a stable limit cycle.

No modifications were introduced into the GER framework. The identical observational pipeline employed in the previous validation stages was executed,

\[
\text{External System}
\rightarrow
\text{Modal Representation}
\rightarrow
\text{Observational Snapshot}
\rightarrow
\text{Persistence Observatory}
\rightarrow
\text{Observational Trajectory}
\rightarrow
\text{Fundamental Geometric Operators}
\rightarrow
\text{Geometric Signature}
\rightarrow
\text{Structural Certificate}
\rightarrow
\text{Observational Partition}.
\]

The Modal Representation immediately revealed a substantially richer spectral structure than those observed for both the Harmonic and Damped oscillators. The Participation Ratio increased from approximately two effective modes to more than five, while the spectral entropy increased by more than a factor of three, indicating a significantly broader modal occupation induced by the nonlinear dynamics.

These differences propagated consistently through the Persistence Observatory. In particular, the entropy observable increased substantially while the remaining observables remained numerically stable throughout the experiment, demonstrating that the observational framework successfully captured the increased dynamical complexity without requiring any modification of the underlying algorithms.

The resulting Geometric Signature was

\[
\begin{aligned}
D &= 1.5608913946393885,\\
C &= 0.2914078484866824,\\
R &= 0.03571428571428571,\\
\Delta &= 2.54\times10^{-16}.
\end{aligned}
\]

Comparison with the previous validation stages revealed a distinct geometric profile. The Diameter decreased relative to both linear systems, the Recurrence reached the largest value observed so far, and the Drift returned to a value numerically compatible with zero, consistent with the stabilization of the trajectory on a limit cycle after the transient regime.

The generated Geometric Signature successfully passed the official Structural Certification process,

\[
\begin{aligned}
\text{Signature Integrity} &:\ \mathrm{PASS},\\
\text{Consistency Check} &:\ \mathrm{PASS},
\end{aligned}
\]

demonstrating once again that structural certification depends exclusively on the internal geometric consistency of the generated signature and remains independent of the physical origin of the analyzed dynamical system.

The certified signature was subsequently incorporated into the official GER reference universe and processed by the observational partition framework. Under the recurrence-induced partition, the Van der Pol oscillator generated a singleton observational block, confirming that the partition machinery consistently accommodates nonlinear external systems using exactly the same mathematical infrastructure previously validated for conservative and dissipative dynamics.

The completion of S27-R3 concludes the first stage of the External Validation Protocol. At this point, GER has successfully processed, certified and partitioned three fundamentally different classes of dynamical systems: a conservative linear oscillator, a dissipative linear oscillator and a nonlinear limit-cycle oscillator. Throughout all three validation stages, the complete observational pipeline remained unchanged, providing strong experimental evidence that the GER observational framework operates as a system-independent geometric observer whose analysis depends solely on the structure of the observed trajectory rather than on the equations governing its dynamics.

These experiments establish the methodological foundation for the next stage of the validation protocol, where increasingly complex chaotic systems will be investigated in order to evaluate the behavior of the Fundamental Geometric Operators under bifurcations, sensitive dependence on initial conditions and high-dimensional attractors.

\subsection{S27-R4 --- External Validation Using the Logistic Map}

The fourth stage of the external validation protocol investigated the behavior of GER when applied to a discrete chaotic dynamical system. Unlike the previous experiments, which considered continuous-time differential equations, this stage employed the classical Logistic Map operating in the fully chaotic regime ($r=3.9$).

This experiment therefore represents the first validation of the GER observational framework on a deterministic discrete system exhibiting sensitive dependence on initial conditions.

As in all previous validation stages, no modifications were introduced into the GER framework. The complete observational pipeline remained unchanged,

\[
\text{External System}
\rightarrow
\text{Modal Representation}
\rightarrow
\text{Observational Snapshot}
\rightarrow
\text{Persistence Observatory}
\rightarrow
\text{Observational Trajectory}
\rightarrow
\text{Fundamental Geometric Operators}
\rightarrow
\text{Geometric Signature}
\rightarrow
\text{Structural Certificate}
\rightarrow
\text{Observational Partition}.
\]

The generated Logistic trajectory remained confined within its expected invariant interval while exhibiting the characteristic irregular evolution associated with deterministic chaos.

The Modal Representation revealed a dominant zero-frequency component, reflecting the non-zero invariant mean of the chaotic sequence. Although the modal energy remained correctly conserved, the spectral distribution differed substantially from the continuous systems investigated previously.

The Persistence Observatory produced an unexpected but internally consistent result. Most observables evolved only weakly throughout the observation window, leading to an observational trajectory occupying a very small region of the geometric space.

Consequently, the resulting Geometric Signature was

\[
\begin{aligned}
D &= 0.010716857466728347,\\
C &= 0.003313136757680963,\\
R &= 1.000000000000000000,\\
\Delta &= 6.27\times10^{-15}.
\end{aligned}
\]

Unlike the signatures obtained for the previous continuous systems, the Logistic Map generated an almost degenerate observational trajectory, characterized by extremely small Diameter and Convergence together with maximal Recurrence.

Despite its unusual numerical values, the generated signature successfully passed the official GER Structural Certification,

\[
\begin{aligned}
\text{Signature Integrity} &:\ \mathrm{PASS},\\
\text{Consistency Check} &:\ \mathrm{PASS}.
\end{aligned}
\]

The certified signature was subsequently incorporated into the official GER reference universe and processed by the observational partition framework, where it generated its own singleton observational block.

Rather than indicating a limitation of the Fundamental Geometric Operators, this experiment revealed an important methodological property of the current external validation protocol. The observational snapshots were generated through circular shifts of a stationary chaotic sequence. Since these shifts preserve most global statistical properties of the Logistic trajectory, successive observational states remained nearly invariant, naturally producing an almost stationary trajectory in the observational space.

This observation identifies a limitation of the current snapshot generation strategy for stationary discrete systems rather than a limitation of the GER observational framework itself. The result therefore constitutes a methodological finding that will guide future developments of the external validation protocol.

The completion of S27-R4 extends the external validation of GER to deterministic chaotic maps while preserving complete compatibility with the original observational architecture. At this stage, the framework has successfully processed conservative, dissipative, nonlinear periodic and discrete chaotic systems using exactly the same geometric pipeline, without requiring any modification of the scientific core.

The next stage of the validation protocol introduces the Lorenz attractor, representing the first continuous strange attractor investigated by GER. Unlike the Logistic Map, the Lorenz system combines continuous evolution with fully developed deterministic chaos, providing an ideal benchmark for evaluating the response of the Fundamental Geometric Operators to high-dimensional chaotic dynamics.
